\documentclass[11pt]{article}
\usepackage[margin=0.95in]{geometry}
\usepackage[T1]{fontenc}
\usepackage{lmodern,microtype,amsmath,amssymb,amsthm,booktabs,array,longtable}
\usepackage{xcolor,hyperref,xurl}
\hypersetup{pdftitle={Gap encodings, congruence finalizers, and a direct proof-stack source for Con(ZFC)},pdfauthor={Research contribution for ProveIt},colorlinks=true,linkcolor=blue!45!black,urlcolor=blue!45!black,citecolor=blue!45!black}
\setlength{\emergencystretch}{3em}
\newtheorem{theorem}{Theorem}[section]
\newtheorem{lemma}[theorem]{Lemma}
\newtheorem{proposition}[theorem]{Proposition}
\newcommand{\ZFC}{\mathrm{ZFC}}
\newcommand{\Con}{\operatorname{Con}}
\newcommand{\Zp}{\mathbb Z_{>0}}
\newcommand{\N}{\mathbb N}
\newcommand{\Z}{\mathbb Z}
\newcommand{\Sq}{\operatorname{Sq}}
\title{Smaller consistency equations through degree separation, congruence norms, and direct proof certificates}
\author{Research contribution for the ProveIt repository}
\date{23 October 2026}
\begin{document}
\maketitle
\begin{abstract}
We compare the six newly supplied constructions and combine positive-gap source coordinates, degree-separated coefficient packing, calibrated carry deletion and a congruence-separated two-root norm. The resulting fixed integer polynomials improve the operation/degree tradeoffs for nine positive and ten unrestricted integer unknowns. The full literal circuits and coefficient payloads accompany the paper. A distinct proof-certificate source avoids the K215 machine and Jones's universal expression; its local polynomial relation is separated carefully from the still-uninstantiated global trace compiler. All consistency equivalences claimed for completed equations are ordinary mathematical proofs formalizable in ZFC, with identified inherited compiler and number-theoretic dependencies. We claim neither global minimality nor a new end-to-end proof-assistant certificate.
\end{abstract}
\setcounter{tocdepth}{1}
\tableofcontents
\clearpage
\section{Results and interpretation}
The requested target is a fixed ordinary integer polynomial $P$, on a stated domain $D$, for which
\[
 \ZFC\vdash \bigl(\neg\exists\mathbf x\in D^n\ P(\mathbf x)=0\bigr)
 \longleftrightarrow\Con(\ZFC).
\]
Here consistency is syntactic: there is no finite derivation of a contradiction in the fixed effective proof calculus for ZFC. We prove an existence equivalence to such a refutation, then negate it. No argument assumes that ZFC is consistent, and no part of the work decides that question.

\begin{center}
\begin{tabular}{llrr}
\toprule
Domain & New construction & Operations & Exact degree\\
\midrule
$\Zp^9$ & Gap $N-q$, simple majorant & 444 & 1,845,940\\
$\Zp^9$ & Gap $N$, simple majorant & 449 & 1,834,516\\
$\Zp^9$ & Reflected gap, simple majorant & 451 & 1,825,540\\
$\Zp^9$ & Gap $N-q$, refined majorant & 476 & 1,520,178\\
$\Zp^9$ & Gap $N$, refined majorant & 481 & 1,510,770\\
$\Zp^9$ & Reflected gap, refined majorant & 483 & 1,503,378\\
$\Z^{10}$ & Reflected gap, congruence norm & 494 & 3,507,876\\
$\Z^{10}$ & Gap $N-q$, congruence norm & 490 & 3,547,076\\
$\Zp^{15}$ & Integral-root transport & 88 & 24\\
$\Zp^{15}$ & Integral-root transport & 87 & 26\\
$\Zp^{15}$ & Integral-root identity & 78 & 30\\
\bottomrule
\end{tabular}
\end{center}
The nine-variable family additionally has explicit nonnegative-domain versions, each with seven extra operations and the same degree. Against the newest incoming nine-positive circuit (473 operations, degree 1,910,404), the 444-operation variant improves both objectives; the refined variants trade some gates for lower degree. Against the incoming ten-integer degree/operation records considered here, both completed ten-integer variants improve substantially. The ten-integer evaluator materializes a fixed threshold constant: retaining the former five constant-only gates would give 499 and 495 operations respectively. This convention does not change either polynomial or its degree.


Every named intermediate register in the completed equations is a polynomial abbreviation, not another quantified variable. Every exponent is a fixed nonnegative integer. There are no divisions, variable powers, square roots, bit operations or unbounded quantifiers inside the final literal circuits. Divisions, radicals and population counts occur only in proofs. The four compiler ports are bound to complete, coherent numerical tuples. Derived large coefficients are materialized as signed binary payloads with hashes rather than left as unspecified universal-polynomial parameters.

\paragraph{What simplicity means here.}
Number of unknowns, their domain, total polynomial degree, arithmetic circuit cost, and coefficient storage are different objectives. We count a live binary addition, subtraction or multiplication as one operation; multiplication by a numeral is charged, and fixed powers are compiled to multiplication gates. Reusing a previously computed node is free. Fixed integer coefficient construction is not included in this circuit cost. Thus a smaller operation count is not a bit-complexity or running-time claim. We retain the relevant tradeoffs rather than declare a single globally simplest polynomial.

\paragraph{A route beyond the previous architecture.}
The new direct proof source does not use the K215 transition table or the Jones universal expression. It addresses the user's broader permission to choose another equiconsistent formalism. Its local transition polynomial and accepting-path equivalence are explicit. The report does not misidentify that local polynomial with a completed consistency equation: converting all finite paths to one fixed ordinary polynomial still has a separately stated implementation boundary.

\section{Recent commits and the incoming frontier}
The review is pinned to ProveIt main
\path{beb0192cb8dc36f78c369c080f2adfbc1887a486}, with the six original archive bytes pinned at arrival
\path{7331974b696693a8eed2aa5f0f64f1c19d9569f0}.
They are placed as contributions 54--59 under
\url{https://github.com/VladimirReshetnikov/ProveIt/tree/beb0192cb8dc36f78c369c080f2adfbc1887a486/Computability/DiophantineConsistency/con-zfc-diophantine-forms}.
The original ZIP identities, retained member hashes and direct-file provenance are in the artifact manifests.

\begin{center}\small
\begin{tabular}{@{}p{4.0cm}lrr@{}}
\toprule
Incoming contribution & Domain & Operations & Degree\\\midrule
54, calibrated carry & $\Zp^9$ & 564 & 4,400,764\\
54, calibrated carry & $\Z^{10}$ & 554 & 14,699,152\\
55, congruence norms & $\Z^{10}$ & 586 & 9,152,284\\
56, free coordinate & $\Zp^9$ & 560 & 4,686,812\\
57, separated degrees & $\Zp^9$ & 499 & 2,628,944\\
57, separated degrees & $\Z^{10}$ & 544 & 8,317,400\\
57, two norms & $\Zp^{12}$ & 89 & 258\\
57, balanced quotient & $\Zp^{12}$ & 94 & 242\\
58, integral root & $\Zp^{15}$ & 79 & 30\\
58, integral root & $\Zp^{15}$ & 88 & 26\\
58, integral root & $\Zp^{15}$ & 90 & 24\\
59, degree separated & $\Zp^9$ & 473 & 1,910,404\\
59, degree separated & $\Z^{10}$ & 533 & 8,469,528\\\bottomrule
\end{tabular}
\end{center}
The table reports the incoming exact-degree claims, except where the independent source audit supplies a fresh degree proof for the 473-operation nine-positive equation. Intake independently checked the literal gate lists, not every incoming semantic theorem. The broader frontier also retains the inherited fifteen-positive 66-operation/degree-157 equation, 75/44, and source 56's 73/54 and 74/46. The new short variants do not dominate all of those points. The two twelve-positive points above remain relevant and are not surpassed by merely comparing with the older degree-897 equation.

\paragraph{Ideas combined, not merely recopied.}
Degree-separated kernels allow linear and quadratic labels to collide safely because they multiply different powers of the packed word. Positive coordinate gaps delete additional coefficient records. A bit whose value is already forced by a recovered congruence can be removed from a carry mask, simultaneously repairing the residue of a leading-block-free code. Finally, coprimality and a fixed congruence distinguish two square classes without a large separator, replacing an eight-factor norm by a quartic norm. The nine-positive branch instead combines the free carrier with the previous sharper quotient majorant and a global addition-chain plan for all base powers.

\paragraph{Formal status.}
Commit \path{dabf4e39f1039a7f77056517ecbb01d6db1a8ced} advances the formal proof of source 52: its literal source gates, residuals, shifted records, selective bounds, low-mask decoding, cofinal bounds, ratio substitution and margin-first finalizer have formal interfaces. The pinned status still distinguishes these pieces from an assembled new source adapter with actual native-port calibration and an internal object-ZFC equivalence. This contribution does not claim that those formal results automatically certify its changed source, carry code or final circuit. Five directly pinned status/proof documents are retained for the precise boundary.

\section{Gap coordinates reduce the quadratic source}
The incoming degree-separated source has $31$ positive coordinates and
$22$ quadratic residuals, with $89$ coefficient records after its two
invertible row operations. A positive-coordinate change reduces this to
$86$ records without changing the finite positional certificate.

Write $q=1+\beta J$, $k=\eta+\zeta$ and $Y=a-E$.
The source equations give
\[
 Y=sN>0,\qquad E=XY>0,\qquad c-\eta=kY>0,
 \qquad P-k=E(c-\eta)>0.
\]
Therefore replacing old $P$ by $T+k$ is a bijection on positive zero sets.
The affected residuals are
\[
 T-Ec+E\eta,\qquad \tau^2-T^2-\eta T-\zeta T-1.
\]
The first loses two records and the second keeps its support. A further
one-record reduction is obtained either by replacing old $N$ by $n+1$,
or by replacing old $L$ by $\ell-1$. These are alternatives, not simultaneous
changes. Their inverse maps preserve positivity because $N=q^3\ge8$ and
$\ell=q^2\ge4$, respectively.

The same $116$-maximum-weight certificate passes all interactions:
$22\,481$ for the $N$-gap variant and $22\,451$ for the $L$-successor variant.
The kernel degrees become $(1719,2020,1929)$, lowering the linear degree
from $2027$; convolution height remains $2161$. The quadratic leading
coefficient stays $+1$. Thus the inherited unconditional positivity estimate
is strengthened, not weakened. Explicit source polynomials, kernels, and an
independent symbolic reconstruction are included.

An optional source lemma reduces the positivity guards from twelve to eight:
$F,\alpha,z,h,s,w,\eta,\zeta$. With native $M_C,M_F\ge0$, transport first
forces $C=q-F-\alpha>0$, and
\[
 R=(q^2-1)((q-1)C+q\alpha)+(M_C+qM_F)J>0.
\]
Hence $j\ge1$ from $R=jc-V$. The identity
$k=jc-V+1+hE$ and $c=kY+\eta>k$ exclude $i=0$ and $g=0$:
in either case the remaining equations force $V\le1$, which would imply
$k\ge c$. If $\gamma=0$, then $k>E$ gives $c>X$ and
$D=X+ac<(a+1)c$, contradicting
$D^2=((a+2)^2-1)c^2+1$. The other coordinate positivities follow directly
from their defining equations. This lemma is not advertised as a gate-count
saving: the degree-separated encoder already reuses its full positive word
in the low digit mask.

A further alternative replaces old $N$ by $n+q$ simultaneously with the
$P$-gap change. Its inverse is $n=N-q=q(q^2-1)>0$, so this is again a
positive-zero-set bijection and leaves only $85$ coefficient records.
Unlike the $86$-record alternatives, this source needs a new positional
layout: the former layout has exactly one unintended target interaction.
Only a separately audited replacement certificate may be used with this
$85$-record version.

\section{Repacked homogeneous layers and positive leading coefficients}
\label{sec:carry-repacking}
The source coordinate changes admit three frozen positional certificates.
The following table records source-code heights, not final equation degrees.
\begin{center}
\begin{tabular}{lrrrrr}
\hline
Certificate & Records & $W$ & $Q$ & $(d_0,d_1,d_2)$ & $H$\\
\hline
Gap $N$ &86&116&2036&$(1571,1927,1836)$&2068\\
Gap $N-q$ &85&116&2050&$(1627,1941,1850)$&2082\\
Reflected gap $N$ &86&116&2031&$(1935,1941,1825)$&2057\\
\hline
\end{tabular}
\end{center}
Here $d_r=\deg D_r$, $Q=\max q_j$, and
$H=\max_{0\le r\le2}(d_r+rW)$.  The incoming split116 height was 2161.
These are certificate-backed improvements, not claims of optimal packing.

For a source monomial $x^\alpha$ of degree $r$, let
$e(\alpha)=\sum_i\alpha_iw_i$.  A coefficient $a_{j,\alpha}$ in row $j$
receives weight $2a_{j,\alpha}$ for degrees zero and one or a diagonal
square, and $a_{j,\alpha}$ for a quadratic cross term.  The reversed kernels
and positive-coordinate word are
\[
 D_r(T)=\sum_{j,|\alpha|=r}c_{j,\alpha}T^{q_j-e(\alpha)},\qquad
 U(T,x)=\sum_i x_iT^{w_i}.
\]
The independent verifier reconstructs all 22 target polynomials in all 31
witnesses and all four independent coefficient ports.  Its finite isolation
check establishes
\[
 [T^{q_j}]\bigl(D_2(T)U(T,x)^2+D_1(T)U(T,x)+D_0(T)\bigr)=2\rho_j(x).
\]
There are 22,481 interactions for either 86-record source and 22,450 for the
85-record source.  The latter needs the new targets: its former incoming
layout had an extra cross-row collision and is not reused here.

\paragraph{A simpler positivity argument.}
The reflected certificate replaces every old position $w_i$ by $117-w_i$
and negates source row 2.  All positions remain distinct and positive; a row
sign change preserves the common zero set.  Its leading coefficients are
\[
 \operatorname{lc}(D_0)=2K_0,\qquad
 \operatorname{lc}(D_1)=2,\qquad
 \operatorname{lc}(D_2)=\beta.
\]
They are all positive integers. If $M$ strictly bounds every absolute kernel
coefficient and $B>M+1$, a polynomial of degree $d$ with positive integral
leading coefficient is greater than
\[
 B^d-M\frac{B^d-1}{B-1}>0.
\]
Thus each $D_r(B)>0$ separately and
$G(B,U)=D_2(B)U^2+D_1(B)U+D_0(B)>0$ for every $U\ge0$.
Unlike the earlier proof, this does not require the quadratic layer to have
strictly greater numerical height than the linear layer.

There is a consequential leading-term tie:
\[
 d_2+2W=d_1+W=2057.
\]
For $U=\sum_iB^{w_i}+g$, the leading $B$-coefficient of $G$ is
$\beta+2$, \emph{not} $\beta$.  The exact final degree certificates account
for this positive sum.

\paragraph{Bounds and completeness.}
The coherent original31 tuple admits the conservative bound
$M=2^{2134927}$ and $B=4096Mb^2$.  The coefficient audit must be repeated
for whichever complete tuple is actually bound; ports are not mixed.
Writing $\ell=\sum_iB^{w_i}$, distinct positions give $\ell<2B^W$.
On $0\le g<4bB^W$, put $U=\ell+g$ and $Z=B^{H+1}$.
A geometric-series coefficient bound gives
\[
 0<G(B,U)<\frac{43Mb^2 B^{H+1}}{B-1},\qquad
 \frac{2G(B,U)}Z<\frac1{32}.
\]
The reflected target gap $H-Q=26$ gives, for $A=\sum_jB^{q_j}$,
$A<Z/64$ and $(B-2)A<Z/64$.  Therefore with
\[
 \begin{split}
 N_0&=B^{W+1},\quad M_0=N_0-1-(b-1)\ell,\quad J_0=2G+A,\\
 S&=g+N_0J_0,\qquad T=M_0+N_0(B-2)A,
 \end{split}
\]
one has $S+T<N_0Z$.  These inequalities concern arbitrary bounded $g$,
before digit recovery.  Once the low mask yields
$g=\sum_i n_iB^{w_i}$, $0\le n_i<b$, the word $U$ encodes $x_i=n_i+1$.
All convolution coefficients have absolute value at most
$993Mb^2<B/4$, so the inherited sparse $+1$ signed-borrow test applies.

Every sufficiently large admissible base represents a fixed positive source
zero.  Row 5 still has value 1 at the all-one tuple, so some digit $n_i$ is
nonzero.  Since all positions are positive, its encoded word is at least $B$
and is divisible by $B$.  These facts preserve the lower-bound and divisibility
conditions used in the integer composition.

The unreflected certificates instead have quadratic dominance gaps
$(25,497)$ and$(25,455)$ and positive quadratic leading coefficient 1; the
usual geometric-series comparison proves their positivity.  Read-only
checks, byte-identical kernel rebuilds and coherent-hash tamper controls
accompany all three certificates.  They authenticate finite data and
polynomial identities, not the inherited unbounded compiler theorem or
an end-to-end kernel proof.

\section{Nine positive unknowns: refined free-carrier composition}
\label{sec:refined-free-nine}

The source transformations and positional certificates above give three
layouts of heights $2068,2082,2057$, each with $W=116$.  All fixed coefficients
are the complete original31 tuple; the high mask is already shifted.
The two finalizers below give the following operation/exact-degree pairs.
\begin{center}
\begin{tabular}{lrrr}
Layout&source records&simple finalizer&refined finalizer\\\hline
$H=2082$&85&$444/1\,845\,940$&$476/1\,520\,178$\\
$H=2068$&86&$449/1\,834\,516$&$481/1\,510\,770$\\
$H=2057$&86&$451/1\,825\,540$&$483/1\,503\,378$
\end{tabular}
\end{center}
These six points are nondominated within the delivered family.  The cheapest
and lowest-degree recommendations improve the incoming nine-positive point
$473/1\,910\,404$ in different ways.  There is no global minimality claim.
The natural-domain translations cost seven further operations each and have
the same exact degrees.

\subsection{The literal calibrated source}
The three kernels $D_0,D_1,D_2$, labels $w_i$, and targets $q_j$ are the
literal data in the selected \texttt{nine/data/H*/} directory.  They satisfy
the proved coefficient-extraction identities for the corresponding positive
source.  Put $M=2^{2134927}$, $\kappa=2134939$, and set
{\small
\begin{align*}
 B&=2^\kappa b^2,&\ell&=\sum_i B^{w_i},& A_t&=\sum_jB^{q_j},\\
 U_s&=\ell+g,&G_s&=D_2(B)U_s^2+D_1(B)U_s+D_0(B),\\
 N_0&=B^{W+1},&E_c&=H+W+2,&M_s&=N_0-1-(b-1)\ell,\\
 S_c&=g+N_0(2G_s+A_t),&&T_c=M_s+N_0(B-2)A_t.
\end{align*}
}
Here $H=\max_d(\deg D_d+dW)$ is a fixed height, not a variable.
In the unreflected layouts the positive leading quadratic kernel dominates
the two others.  In the reflected layout all three kernels are independently
positive for $B>M+1$.  Its quadratic and linear heights tie, and their
leading coefficients add to $\beta+2>0$.  Thus $G_s>0$ before a final zero
or a binary-base hypothesis is assumed.

On the branch $g<4bB^W$, the exact geometric estimates give
$2G_s/B^{H+1}<1/32$, while each target contribution is below $1/64$ of
that block.  Hence $S_c+T_c<B^{E_c}$.  The sparse borrow test yields
\[
 S_c\mathbin{\&}T_c=0
 \quad\Longleftrightarrow\quad
 g\text{ encodes a positive common source zero}
\]
when $b$ is a power of two.  This uses the coefficients bounded by
$993Mb^2<B/4$ and the target digit $4\rho_j+1$ tested against $B-2$.

Let $H_p=E_c+W+23$, and define the fixed integers
{\small
\begin{align*}
 \delta&=(H_p-15)\bmod2,&\lambda&=(22-H_p\kappa)\bmod2,\\
 e_N&=(H_p-15-\delta)/2,&e_2&=((15+\delta)\kappa+\lambda-22)/2.
\end{align*}
}
Use the polynomial registers
\[
 L=2^\lambda bB^{E_c},\qquad N=2^{e_2}b^\delta B^{e_N},\qquad
 R=L(S_c+T_c+1)-S_c-1,\qquad t=bB^W.
\]
The exact two-block population identity proves
\[
 N^2\mid {2R\choose R}\quad\Longleftrightarrow\quad S_c\mathbin{\&}T_c=0.
\]
Indeed, if $b=2^k$ and $B=2^v$, then
$\operatorname{pop}(T_c)=(W+23)v-31k-22$ and
$\log_2L+\operatorname{pop}(T_c)=2\log_2N$.
Unconditionally $R,N\ge8$ and $b\le R,N$.  The endpoint $b=1$ is harmless:
its low mask would force $g=0$, contrary to the positive domain.

\subsection{The nine-input equation}
The nine inputs are exactly $b,g,h,s,w,C,i,j,n>0$.  Set $j_0=j-1$ and
{\small
\begin{align*}
 U&=N^2w,&Y&=N^2s,&M_0&=RY,&P&=2M_0^2U,\\
 K&=R+1+h(P-1),&A&=M_0(U+1),&B_1&=2R+1,&W_1&=bw,\\
 J&=4A-5,&Q&=3W_1C-2(W_1^2-1),&D&=(A^2-1)C^2+1,\\
 k_0&=iJC,&E&=k_0DC,&F&=(A^2-1)E^2+1,\\
 \Lambda&=CD+1-A,&G&=A+F\Lambda,&H_1&=C+F(B_1+j_0C),\\
 I&=(G^2-1)H_1^2+1,&a_0&=(P^2-1)K^2+1,&a_1&=DFI,\\
 m&=K^2(4t-g)-16t(C-KY)^2,&\Omega&=Jm,&V_0&=PK,\\
 E_0&=J(n+W_1^2)+Q.&&&&
\end{align*}
}
All symbols are polynomial abbreviations, not extra quantified unknowns.
For the refined version put
{\small
\begin{align*}
 B_8&=8D^2-4D-1,\qquad N_* =2G+\Lambda,\\
 T_6&=k_0\{128B_8N_*-k_0^2[16(B_8+D)-8]\Lambda
                       +C(14k_0^2-192)\},\\
 T&=T_6+i^3\{C[C(576A-2151)-(896A^2-4256A+8442)]+1024\}
          +288i+2048,\\
 V_1&=TH_1,\qquad q_*=2048^2,\qquad
 \Theta=E_0-\Omega(V_0+V_1).
\end{align*}
}
The single equation is
\begin{equation}
 \boxed{\bigl[\Theta^2-\Omega^2(a_0+q_*a_1)\bigr]^2
             -4q_*\Omega^4a_0a_1=0.}
 \label{eq:round9-refined-nine}
\end{equation}
The cheaper simple version replaces $V_1$ by
$k_0(8D^2-4D-1)GH_1$ and $q_*$ by $64$.

\subsection{Why a free carrier preserves the equivalence}
The refined majorant proof retained from the prior construction needs only
$A\ge729$, $C\ge1$, and $i\ge1$.  Here $A\ge33280$, so it applies without
the old carrier substitution and gives
\[
 V_0>\sqrt{a_0},\qquad V_1>2048\sqrt{DFI}>0.
\]
Also
\[
 E_0=Jn+(J-2)W_1^2+3W_1C+2>0.
\]
A zero of \eqref{eq:round9-refined-nine} cannot have $m=0$.
Its real two-root factorization gives
\[
 E_0/J=m(V_0+V_1+r_0+r_1),\qquad
 r_0^2=a_0,\quad r_1^2=q_*a_1.
\]
The strict majorants make the parenthesis positive for every sign choice;
therefore $m>0$ is obtained \emph{before} any integrality assertion.
It follows that $|C-KY|<K/2$, hence $C>2K>2R+1+W_1$.
This recovers the old carrier range.  The usual ratio inequalities now give
$a_1>a_0$, so $r_0+r_1=\Theta/\Omega$ is rational and nonzero.
Squaring isolates both roots, which are rational and consequently integral;
$r_1/2048$ is integral as well.  The displayed factor identity gives $J\mid Q$.
The exact CRT replacement using $F\equiv1\pmod{JC}$ therefore supplies all
four inherited ratio conditions: the two squares, divisibility, and positive
margin.  Jones's theorem and the source decoder yield a syntactic refutation.

Conversely, take the ratio witnesses supplied by a source zero and the
positive roots.  Choose
\[
 n=m(V_0+V_1+r_0+r_1)-W_1^2-Q/J.
\]
It is integral.  Since $A>W_1$, $D\ge A^2$, $D>C^2$, and
$V_1>2048D$, one has $W_1^2+Q/J<2D<V_1$, so $n>0$.
One norm factor vanishes.  This proves the existence equivalence and,
by negation, the ZFC-provable equivalence between no positive zero and
$\operatorname{Con}(\mathrm{ZFC})$.  The inherited compiler and numerical
ratio theorem remain identified dependencies; no new end-to-end kernel
certificate is claimed.

\subsection{Exact degree and implementation checks}
Let $r=\deg R$, $\nu=\deg N$, $c_0=2W+1$, and $a=r+4\nu+2$.
The exact local defect $T^2-2048^2DFG^2$ has703 terms, each with exponents
at most $(15,18,6)$ in $(A,C,i)$.  The unique term attaining all three
maxima is $9\cdot2^{18}i^6A^{15}C^{18}$.  Its complete integer expansion
is independently reconstructed.  Consequently, for
$\Delta_1=V_1^2-q_*a_1$,
\[
 \Delta_1=H_1^2(T^2-q_*DFG^2)+q_*DF(H_1^2-1),\qquad
 \deg\Delta_1=31a+2r+44.
\]
Resolve the main cancellation by the exact identity
{\small
\begin{align*}
 \Theta^2-\Omega^2(a_0+q_*a_1)
 ={}&E_0^2-2E_0\Omega(V_0+V_1)\\
 &+\Omega^2[(K^2-1)+\Delta_1+2V_0V_1].
\end{align*}
}
The term $\Omega^2\Delta_1$ uniquely dominates, and its square dominates
the remaining norm term.  Thus the refined exact degree is
\[
 \boxed{86r+328\nu+4c_0+260},\qquad
 r=4E_c-1,\quad\nu=2e_N+\delta.
\]
Its diagonal leading coefficient is
$81\cdot2^{88}\tau^4\rho^{86}\sigma^{328}>0$, where
$\rho,\sigma,\tau$ are the leading coefficients of $R,N,t$.
The simpler finalizer has degree $104r+400\nu+8W+316$.

The evaluator charges every binary addition, subtraction, and multiplication,
including multiplication by constants.  The new global backward addition
chains share required base powers across all three kernels, the masks and
the threshold.  They do not treat exponentiation as one operation.
The package includes all12 literal circuits, their fixed coefficient payloads,
1080 independent actual-coefficient modular replays, full coefficient-recipe
checks, exact degree inequalities, and the exact703-term identity.

\section{A congruence-separated quartic norm in ten integers}
\label{sec:round9-integer}

The ten-integer route combines the new positive-gap, homogeneous-degree
source with the two-root congruence finalizer of the incoming
\emph{Congruence Norms} contribution. A further one-bit deletion removes
the outer leading block without losing the congruence required by the
arithmetic bridge. The result is an ordinary integer polynomial, not an
exponential Diophantine formula. All named registers below are abbreviations,
not additional unknowns. The exact coefficient payloads and literal circuit
are included in \texttt{integer/}.

\subsection{The leadless source and its calibration}

Use the gap source and its exact three-kernel certificate from the source
section. Write $W=\max_i w_i$, $H=\max_d(\deg D_d+dW)$, and
$A_T=\sum_jB^{q_j}$. Supply exactly
\[
 t,g_0,h_0,k_0,l,w,x,y,z,m\in\mathbb Z.
\]
In the primary mirror layout, row2 is additionally negated, an invertible
row operation; the kernel identity uses these oriented source rows.
Set $\kappa=2134940$, $\lambda=2^\kappa$, and define
{\small
\begin{align*}
 b&=6t+4,\quad g=bg_0,\quad h=bh_0,\quad k=b(3k_0+2),\quad B=\lambda b^2,\\
 \ell&=\sum_i B^{w_i},\quad u=\ell+g,\quad N_0=B^{W+1},\quad N=B^{W+H+2},\\
 G_c&=D_2(B)u^2+D_1(B)u+D_0(B),\quad M_c=N_0-1-(b-1)\ell,\\
 S_c&=g+N_0(2G_c+A_T),\quad T_c=M_c+N_0(B-2)A_T-2.
\end{align*}
}
The source ports for the bridge are
\begin{equation}\label{eq:r9-ten-source}
 \boxed{\quad X=N(S_c+T_c+1)-S_c-1,\qquad
 Y=2^{16\kappa-23}bB^{H+2W+9},\qquad
 \mathcal C=bB^W.\quad}
\end{equation}
No division appears in these expressions. The coefficient $2^{16\kappa-23}$
is materialized, not left as an unevaluated algorithmic parameter.

All signed assignments have $b\ne0$, $2\mid b$, and $b\equiv1\pmod3$.
Because all source weights are positive, direct reduction proves
\[
 X\text{ odd},\qquad b\mid X+1,\qquad b\mid Y,\qquad 3\mid X.
\]
Indeed $B,N,N_0\equiv1\pmod3$, $\ell\equiv31$, $A_T\equiv22$, so the
old high mask is $-22\equiv2$, the new $T_c$ is zero, and $X\equiv T_c$.

The deletion of $2$ from the mask is justified only after the finalizer has
recovered that $b$ is a positive even square. At that point $b\ge4$ and
$4\mid g$. The low base-$B$ digit of the old mask is $B-1$, while bit1 of
$S_c$ is zero. Thus clearing bit1 preserves binary disjointness and reduces
the mask population by exactly one. The usual source bounds give
$S_c+T_c<N$ for $0<g<2\mathcal C$.

If $b=2^s$ and $B=2^v$, then
\[
 \operatorname{pop}(T_c)=(W+23)v-31s-23.
\]
The two disjoint blocks of $X=(S_c+T_c)N+(N-1-S_c)$ give
\[
 \operatorname{pop}(X)=\log_2N+\operatorname{pop}(T_c)
 -\bigl(\operatorname{pop}(S_c)+\operatorname{pop}(T_c)
              -\operatorname{pop}(S_c+T_c)\bigr).
\]
The maximum is $(H+2W+25)v-31s-23=\log_2Y$. Consequently
\[
 Y\mid\binom{2X}{X}\quad\Longleftrightarrow\quad S_c\mathbin{\&}T_c=0,
\]
and the source certificate decodes a positive common source zero. This is
an exactly recalibrated threshold, smaller by four than the old leading-block
threshold. It does not assume $4\mid b$ on arbitrary signed inputs.

\subsection{The complete polynomial}

Define the division-free bridge registers
{\small
\begin{align*}
 L&=lY,&U&=2LX,&V&=4gwY,\\
 W_b&=bw,
 A&=U(V+1),&B_1&=2X+1,&C&=B_1+(A-2)h,\\
 P_{1/2}&=(LX)UV,&P&=2P_{1/2},&K&=X+1+k(P-2),\\
 S&=2A-5,&T&=3W_bC-2(W_b^2-1),&
 Q&=(P_{1/2}^2-1)K^2+1,\\
\end{align*}
}
{\small
\begin{align*}
 \Delta&=8\mathcal C^3gK^2
       -g^2\bigl(32(C-KL)^2\mathcal C^3+g^2K^2\bigr).
\end{align*}
}
The subscript $1/2$ names an explicitly multiplied polynomial; it is not an
instruction to divide an unknown. Put
{\small
\begin{align*}
 \xi&=zx,&D&=(A^2-4)C^2+4,&E&=C^2D\xi,\\
 F&=4(A^2-4)E^2+1,&G&=2(A+2)(A-2)^2E^2-1,\\
 \Omega&=16E,&H_c&=C+(B_1-C)F+\Omega Fy,&I&=(G^2-1)H_c^2+1,\\
 \Sigma&=\xi S^5,&\rho&=T+Sy,&
 \Theta&=(\rho^3-\rho S^2-S^3)(\rho^2+23S^2),\\
 r&=bQ,&a&=r+(3\Delta-1)(rz)^2,&d&=DFI,\\
 Z&=\Sigma m-\Theta.
\end{align*}
}
The final equation is
\begin{equation}\label{eq:r9-ten-final}
 \boxed{F_{10}=(Z^2-\Sigma^2(a+d))^2-4\Sigma^4ad=0.}
\end{equation}
The four fixed compiler ports are inherited as a whole from the original31
instance. The positive-gap transformations change its coordinates, not its
coefficient binding.

\subsection{Why the two-root norm is sufficient}

The normalizations imply $2b\mid U,A$, $b\mid K$, $3\mid K$, and
$Q\equiv1\pmod{b^2}$, $Q\equiv1\pmod3$. Also $2b\mid C+1$ and
$4b\mid D$. Therefore, without any positivity assumption,
\[
 a/b\text{ is odd},\qquad d/b\text{ is even},\qquad a\ne d.
\]
The odd integer $S$ is never zero. The quintic
$\pi(v)=(v^3-v-1)(v^2+23)$ has no rational root, so $\Theta\ne0$.
At a zero of \eqref{eq:r9-ten-final}, $\Sigma\ne0$. With
$q=Z/\Sigma$, the elementary identity
\[
 (q^2-a-d)^2=4ad
\]
and $a\ne d$ force $q\ne0$. The rational numbers
\[
 u_0=\frac{q^2+a-d}{2q},\qquad v_0=q-u_0
\]
satisfy $u_0^2=a$, $v_0^2=d$. They are integers because their squares
are integers. Hence both radicands are squares and $\Sigma\mid\Theta$;
the standard valuation argument gives $S\mid T$.

The coprime factors of
$a=r[1+(3\Delta-1)rz^2]$ are signed squares. Since $r\equiv1\pmod3$,
they must both be positive squares. The same argument for $r=bQ$, using
$\gcd(b,Q)=1$ and $b\equiv Q\equiv1\pmod3$, makes $b,Q$ positive
squares. Then $b\ge4$, and the second factor of $a$ forces $\Delta>0$.
The definition of $\Delta$ gives $K\ne0$ and $0<g<2\mathcal C$.
If $l=0$, then $P=0$ and $Q=1-K^2\le0$, a contradiction. Thus the
nonvanishing condition required by the inherited Sun bridge is recovered
rather than imposed as an additional polynomial factor. That bridge now
gives the central-binomial divisibility and hence a source zero.

For completeness, choose any sufficiently large $b=4^s$ in the source
encoding. Positive source weights give $b^2\mid g$ and
$X+1\equiv-g\pmod{b^2}$. The constructive bridge supplies the exact
Lucas values $C=u_{2X+1}(A)$ and $K=u_{X+1}(P)$. The recurrence and the
facts $X$ odd, $4b^2\mid P,A^2$ show that its existing witnesses satisfy
$b\mid h$, $b\mid k$, and $k/b\equiv2\pmod3$. Thus the displayed
restrictions lose no canonical solutions. Pell's theorem for the positive
nonsquare $(3\Delta-1)r$ chooses $z$ making $a$ square. The inherited
residue interpolation chooses $\xi$ divisible by $z$ and a suitable $y$
while keeping $d$ square and making $\Sigma\mid\Theta$. Then
$m=\Theta/\Sigma+\sqrt a+\sqrt d$ makes \eqref{eq:r9-ten-final} vanish.

These arguments compose with the source/refutation theorem to establish,
in ZFC, that the absence of integral solutions is equivalent to
$\operatorname{Con}(\mathrm{ZFC})$. The complete soundness and completeness
proof, including every source bound and congruence-class reconstruction,
is in \texttt{integer/PROOF.md}. The bridge references are Sun's
2021 Theorems4.12--4.13 and the modified certificate/interpolation results
in his 2026 ten-unknowns preprint. They remain explicit mathematical
dependencies, not claims of newly checked end-to-end kernel verification.

\subsection{Exact degree and evaluator accounting}

Write $r_X=4(W+H)+6$, $r_Y=2H+4W+19$, and $a_A=r_X+2r_Y+4$.
The dominant final term is uniquely $\Sigma^4d^2$, and
\[
 \deg d=88a_A+126,\qquad \deg\Sigma=5a_A+2.
\]
Thus the exact total degree is
\[
 \boxed{\deg F_{10}=196a_A+260=2352W+1568H+9668.}
\]
Full literal circuits, live-gate counts, actual-tuple modular replay, and
nonzero highest-degree coefficients over three primes independently check
the stated metrics. Global shared power chains reduce evaluation gates
without changing a single polynomial coefficient. The fixed threshold is
materialized as one leaf; charging the five constant-only gates in the older
evaluator convention adds five operations, with no effect on degree.

\section{Smaller evaluators for the integral-root equations}
The incoming integral-root mechanism is retained, rather than advertised as a new discovery here. Three literal fifteen-positive circuits improve from operation/degree pairs $90/24$, $88/26$ and $79/30$ to
\[
 (88,24),\qquad(87,26),\qquad(78,30).
\]
All retain the entire incoming $590069$-bit-radix tuple, not the original31 tuple used by the nine- and ten-variable branches.

\subsection{The explicit degree-24 polynomial}
The positive inputs are
\[
 j,F,\alpha,z,f,h,i,T,X,Y,r,\eta,\zeta,u,\gamma.
\]
Define the following polynomial abbreviations:
\begin{align*}
 J&=3j+1,&q&=\beta J+1,& C_w&=q-F-\alpha,\\
 E&=XY,&k&=\eta+\zeta,&c&=kY+\eta,\\
 a&=E+Y,&H&=4a+3,&\Delta&=a^2+H,\\
 D_0&=X+ac+\gamma H,&\ell&=EkY,&I&=ic^2,\\
 R&=((q-1)(q-F)+\alpha)(q^2-1)+(M_C+qM_F)J,\\
 \widehat R&=k-hE-1,&U&=\Delta(f^2-1),&V&=f^2(cT-\widehat R)-c.
\end{align*}
The residuals are
\begin{align*}
 A_0&=[q^3(\ell+r)+Y]^2-q^6[\ell(\ell+k)+1],\\
 A_1&=(qD_0+XH)^2-q^2(\Delta c^2+1),\\
 A_k&=\widehat R-R,\\
 A_t&=q[K_0C_w-F-z(q-1)]+XC_w,\\
 A_a&=(V+u)^2-Uu\bigl(V+(V+u)\bigr)-1,\\
 A_s&=I^2-U.
\end{align*}
Then the completed equation is
\[
 \boxed{A_0^2+A_1^2+A_k^2+A_t^2+A_a^2+A_s^2=0.}
\]
Its literal evaluator, with common subexpressions shared, has 88 live binary gates. The highest homogeneous part remains
\[
 4X^4Y^6f^8u^2\bigl((\eta+\zeta)T+hX\bigr)^2,
\]
which is nonzero and has degree 24.

\subsection{Why both changes preserve the theorem}
First share the already computed auxiliary gap: $2V+u=V+(V+u)$. This is an identity over every commutative ring and saves one addition in each of the three circuits.

For the degree-24 case only, the old transport residual is
\[
 A_t^{\rm old}=q[K_0C_w+(q-F)-z_{\rm old}(q-1)]+XC_w-q.
\]
Our residual obeys the exact all-input identity $A_t(z)=A_t^{\rm old}(z+1)$, while no other residual uses this transport coordinate. Any positive new zero is therefore an old positive zero with $z_{\rm old}=z+1$. The converse uses canonical witnesses, not a claim about arbitrary old zeros. The retained ordinary compiler proof gives
\[
 w=2^R/q>q>F,\quad C_w>0,\quad
 (K_0+w)C_w-F=(z_{\rm old}-1)(q-1).
\]
The left side is positive, and $q>1$, so the integral $z_{\rm old}$ is at least two. Thus $z=z_{\rm old}-1$ is a valid positive new coordinate. All subsequent integral-root reconstructions leave that transport coordinate unchanged. This supplies completeness.

The old integral-root theorem recovers $q^3\mid Y$ and $q\mid X$ by rational-root integrality in the first and main monic square equations. Its congruence $J=3j+1$ is retained: deleting it without an alternative proof of $\gcd(q,3)=1$ is not justified. The inherited canonical Pell gaps and signed rank argument then give the refutation equivalence. The full predecessor proof and complete coefficient binding are retained as inert sources. Translation by a constant preserves degree, and the other two modified outputs are polynomial-identical to their predecessors.

The authored checker validates all source hashes, literal gates and the two identities. It performs 360 exact signed and 360 modular whole-circuit regressions, plus 6000 local checks. These are implementation checks; the two universal identities and canonical positivity argument above are the proof of the changes.

\section{A genuinely different source: one-stack proof certificates}
\label{sec:onestack}
The user explicitly permits abandoning both K215 and Jones's universal
expression. The following construction does so: it recognizes a proof
nondeterministically, rather than arithmetizing a deterministic enumerator.
Its finite accepting paths are mathematically equivalent to a ZFC refutation.
The local polynomial is fully materialized. A closed global compiler-input
proposal is supplied, but has not been elaborated, kernel compiled, or
expanded to count its final variables. In particular, the local equation
alone is \emph{not} a consistency equation.

\subsection{Numerical calculus and stack}
Use the numerical calculus $C_{222}$ of the pinned repository, not its machine
implementation. Its independent semantic calibration is
\[
 C_{222}(1)\ \longleftrightarrow\ \text{there is a ZFC refutation}.
\]
The precise source lemmas are
\path{KEndpoint.c222_iff_zf_refutation} and
\path{ConstructibleChoice.zf_refutation_iff_zfc_refutation}.
Put $I(a,b)=2(a+b)(a+b+1)+4b$, $A(a,b)=I(a,b)+2$,
$M(a,b)=I(a,b)+1$, $E(a,b)=I(a,b)+3$, $N(a)=I(a,1)$, and
$X(a,b)=N(A(a,N(b)))$.
The eight unconditional axiom codes are
\begin{align*}
 r_1&=X(a,E(a,b)),\\
 r_2&=I(A(a,I(b,c)),I(A(a,b),A(a,c))),\\
 r_3&=I(E(a,b),I(M(c,a),M(c,b))),\\
 r_4&=I(A(a,A(b,c)),A(b,A(a,c))),\\
 r_5&=I(E(c,a),I(M(c,b),M(a,b))),\\
 r_6&=I(X(a,A(a,b)),b),\\
 r_7&=I(E(c,a),I(E(c,b),A(M(a,b),E(a,b)))),\\
 r_8&=I(I(I(c,u),a),I(I(a,c),I(b,c))).
\end{align*}
The remaining guarded code $B(a,b,c)$ is defined by
\begin{align*}
 R(b,k,m)&=A(b,I(m,N(A(b,I(m,N(M(b,k))))))),\\
 C(a,b,k,m)&=A(b,I(M(b,k),N(A(a,I(m,\\
 &\hspace{18mm}N(A(a,I(A(k,I(M(k,a),M(k,b))),R(b,k,m))))))))),\\
 B(a,b,c)&=N(I(I(M(a,b),A(c,M(a,b))),\\
 &\hspace{18mm}A(c,I(C(a,b,c,c),N(R(b,c,c)))))).
\end{align*}
It is available when $c\ne a,b$. All formulas are also included as literal
binary arithmetic DAGs in \texttt{stack\_step.json}.

Encode a list by $\operatorname{Cons}(h,t)=1+(h+t)(h+t+1)/2+t$,
with empty list $0$. Every positive natural number has a unique head and
tail, and the tail is strictly smaller. Hence every natural number is a
finite stack code. Define the polynomial
\[
 L(s,h,t)=2s-2-(h+t)(h+t+1)-2t.
\]
Over natural numbers, $L(s,h,t)=0$ is exactly $s=\operatorname{Cons}(h,t)$.
The accepted stack is $\operatorname{Cons}(1,0)=2$.

\subsection{Three branches, five shared witnesses}
With natural witnesses $a,b,c,u,v$, put
\begin{align*}
 G&=(c-a)^2(c-b)^2-(u+1)^2,\\
 H&=v^2\prod_{j=1}^{8}(v-r_j)^2\bigl((v-B(a,b,c))^2+G^2\bigr),\\
 T_{\rm ax}&=L(s',v,s)^2+H,\\
 T_{\rm mp}&=L(s,I(a,b),v)^2+L(v,a,c)^2+L(s',b,c)^2,\\
 T_{\rm gen}&=L(s,b,c)^2+L(s',A(a,b),c)^2,\\
 T&=T_{\rm ax}T_{\rm mp}T_{\rm gen}.
\end{align*}
All three factors are nonnegative. Thus $T=0$ means exactly one of the
following permitted transitions (possibly more than one): push truth or an
axiom; replace adjacent top entries $I(a,b),a$ by $b$; replace top entry $b$
by $A(a,b)$. For the guarded axiom, $G=0$ implies $c\ne a,b$ because $u+1>0$.
Conversely, for $c\ne a,b$ choose $u=|(c-a)(c-b)|-1$.
The same $u$ can serve as the fourth parameter of $r_8$ in a different
branch: an existential disjunction does not require distinct witness slots.

\paragraph{Exact correctness of finite paths.}
Every stack reached from $0$ contains only $C_{222}$ theorems, by induction
on the path length and the three inference rules. Consequently reaching
$2$ implies $C_{222}(1)$. Conversely, prove by induction on a derivation
that its conclusion $h$ can be pushed onto \emph{any} pre-existing stack
$S$, without changing its tail. An axiom is pushed directly. For modus
ponens first execute the derivation of $a$, then that of $I(a,b)$, then
apply the pop-two/push-one branch. For generalization execute the premise,
then apply the third branch. Thus a derivation of $1$ yields a path from
$0$ to $2$. Duplicating a premise merely increases time; no history register,
cursor, or saved proof list is required. This proof is finitary and
formalizable in ZFC, so the inherited semantic calibration gives
\[
 \exists\text{ a finite }T\text{-path from }0\text{ to }2
 \quad\longleftrightarrow\quad \neg\operatorname{Con}(\mathrm{ZFC}).
\]

\subsection{Materialized local quartic and closed global proposal}
The shared binary DAG has 382 operation nodes, seven inputs, and numeral
leaves at most four. Its syntactic degree upper bound is 1,048,802.
Inline affine arithmetic and introduce a nonnegative wire only at each
nonlinear product. Every such product is a square or a product of
nonnegative expressions; this is why the guard uses separate squares.
There are 84 wire variables. Together with $a,b,c,u,v$, the local source
has 89 natural witnesses. Its 85 quadratic residuals have a sum of squares
of exact degree four with 3,108 monomials. There are 91 total natural variables: two exposed stack states and 89 existential witnesses. These are \emph{local} counts,
not counts for the final consistency equation.

For the global closure take only three outside variables $t,x,y$, and
$\beta(x,y,i)=x\bmod(1+(i+1)y)$. Require
\[
 \beta(x,y,0)=0,\quad \beta(x,y,t)=2,\quad
 \forall i<t\ \exists\mathbf w\in\mathbb N^{89}\quad
 Q(\beta(x,y,i),\beta(x,y,i+1),\mathbf w)=0.
\]
The CRT coding lemma supplies $x,y$ for every finite stack sequence.
Conversely the displayed conditions decode a valid finite path.
No extra $t>0$ condition is needed: $t=0$ contradicts the two endpoints.
The displayed bounded formula is not itself an ordinary Diophantine
polynomial. The artifact \texttt{OneStackEquation.lean} specifies its
closed constructor translation: remainder graph representations, finite
composition, \texttt{Rep.boundedForall}, existential closure of $(t,x,y)$,
and finally \texttt{QuarticCompiler.compile}. This bypasses Jones entirely.
Its primitive constructors are explicit data-valued representations in the
pinned repository, not an unspecified appeal to MRDP.

The local result and path equivalence above are mathematical proofs. The
new Lean constructor proposal remains uncompiled, and its complete output
is unmaterialized and uncounted. It is therefore a promising alternate
architecture, \emph{not a claimed smaller completed consistency equation}.

\section{Reproduction, provenance and limitations}
The ZIP contains the editable paper, compiled PDF, exact arithmetic circuits, explicit signed coefficient payloads, source and packing certificates, retained proofs and licenses, and newly authored verification programs. Search executables and unrelated incoming programs are not required to evaluate or check a delivered equation. The incoming archives were inspected as inert data; no upstream Python module was imported or executed.

\subsection{Three different levels of evidence}
\begin{enumerate}
\item \emph{Byte provenance.} The original six archives have exact Git blob and SHA-256 identities. Retained members are matched to those archives, and directly retrieved formal sources are checked by their own Git blob identity. The sealed release manifest checks every included file.
\item \emph{Finite arithmetic certificates.} These reconstruct source coordinate substitutions, check every relevant degree-specific target interaction, rebuild literal circuits, verify complete fixed-port bindings, and use exact identities and actual-coefficient leading-term witnesses. Modular samples are regression tests, not the sole basis for a polynomial identity or an upper degree bound.
\item \emph{Universal mathematics.} Positivity recovery, carry normalization, rational square-root integrality, square-class separation, CRT and canonical Pell reconstruction prove the new existence equivalences. Their identified compiler/calculus and number-theoretic dependencies remain inherited. No finite search establishes consistency.
\end{enumerate}

The finite checks cover all 22 target polynomials over 31 witnesses and four independent coefficient ports, not just sample values of the huge constants. Source packing is checked at 22,481 interactions for the 86-record layouts and 22,450 for the 85-record layout. The integral-root simplifications include 720 full-circuit exact/modular regression comparisons and 6,000 local identity checks. The refined majorant has a separately reconstructed exact 703-term defect polynomial. The direct proof source is checked against its literal local polynomial and quartic expansion. Deliberately corrupted packing coefficients and target positions are rejected even after their code hashes are updated coherently; these tests are distinct from sealed-file hash rejection.


\subsection{How to replay}
From a clean extracted archive, run \texttt{python verify.py}. The driver first checks the sealed manifest and source provenance, makes an isolated copy, and runs only the enumerated authored checkers and builders. Its rebuild ledger compares generated artifacts byte-for-byte with the sealed outputs. Run \texttt{python negative\_controls.py} for deliberate corruptions; a rejected byte tamper is recorded separately from a rejected mathematical certificate. Build the PDF with \texttt{bash report/build.sh}. PDF metadata can differ across TeX environments even when the mathematical text is unchanged.

The two small-variable branches intentionally use the complete original31 coefficient tuple; the short integral-root branch retains its different complete $590069$ tuple. No individual port is borrowed from a smaller tuple. The coefficients are compressed losslessly for storage; the circuit definition uses their integer values. These storage choices are not a new coefficient-height record.

\subsection{Claims deliberately not made}
The report neither proves nor refutes $\Con(\ZFC)$, establishes a global minimum, changes positive unknowns into unrestricted integer unknowns without an explicit transformation, nor presents the new full equations as kernel-checked. The direct proof-certificate local polynomial is not itself a consistency equation. The reported arithmetic-operation count excludes fixed numeral construction and says nothing about measured evaluation time with enormous integer values.

\section{Further research questions}
\begin{enumerate}
\setlength{\itemsep}{0pt}\setlength{\parskip}{0pt}
\item Instantiate and measure the direct proof-certificate trace compiler. A small local relation is not automatically a small final Diophantine equation; bounded-universal elimination and witness coding must be included in the count.
\item Optimize the arithmetic proof calculus itself, independently of a machine implementation. Junk formula codes are permissible if the numerical closure relation retains the syntactic refutation equivalence.
\item Search jointly over positive-coordinate changes, reversible row operations, degree-specific labels and target positions. Optimizing a positional code before simplifying the source can freeze avoidable records into the objective.
\item Develop a proof-producing multi-output addition-chain planner for the common radix powers. The current circuits are explicit and replayable, but the search is not a proof of global circuit optimality.
\item Determine whether the congruence-separated two-root finalizer can replace other high-degree square-class mergers while preserving unconditional signed soundness and canonical completeness.
\item Transfer additional divisibility conditions into existing monic roots without increasing their cleared-denominator degrees. The failed naive changes in the integral-root branch show why denominator cost and positive witness reconstruction must be analyzed together.
\item Formalize the changed source, carry and finalizer adapters against the existing Lean/Rocq endpoint. Exact finite certificates and modular degree witnesses do not replace the universal semantic proof.
\end{enumerate}
\paragraph{Conclusion.} The completed equations improve explicit tradeoffs. The separate direct proof source replaces enumeration by existential proof certificates, but its global trace elimination remains unmaterialized.

\end{document}
